% ══════════════════════════════════════════════════════════════
% Chapter 3 — Risk Systems and Value-at-Risk
% ══════════════════════════════════════════════════════════════

\chapter{Risk Systems and Value-at-Risk}
\label{ch:risk-systems}

\begin{projectconnection}[Project Connection --- Phase 0 \& Phase 2]
This chapter teaches you what every SecDB risk output actually measures.
In Phase~0 you will audit which risk cube nodes are accessible; this chapter
ensures you understand the data before engineering features.
In Phase~2, every feature family---VaR dynamics, factor concentration,
scenario P\&L, VaR utilization---maps directly to concepts here.
When you build features in Chapter~\ref{ch:features}, you will reference
the mathematics and intuitions developed below.
\end{projectconnection}


% ──────────────────────────────────────────────────────────────
\section{What VaR Measures}
\label{sec:var-definition}
% ──────────────────────────────────────────────────────────────

\begin{prereq}[Quantiles and the Normal Distribution]
A \textbf{quantile} $q_\alpha$ of a random variable $X$ is the value below
which a fraction $\alpha$ of the probability mass lies:
$\Pr(X \le q_\alpha) = \alpha$.
For a normal distribution $\N(\mu,\sigma^2)$, the $\alpha$-quantile is
$q_\alpha = \mu + \sigma\,z_\alpha$, where $z_\alpha$ is the standard
normal quantile.  Key values: $z_{0.95}=1.645$, $z_{0.99}=2.326$,
$z_{0.975}=1.960$.
\end{prereq}

Value-at-Risk is the single most important number produced by a bank's risk
system.  Every trading desk has a VaR limit, every risk report leads with VaR,
and every regulator demands it.  The concept is simple: how much could we lose
on a bad day?

\begin{definition}[Value-at-Risk]
Let $L$ denote the portfolio loss over a holding period $h$ (typically 1 day
or 10 days).  The \textbf{Value-at-Risk} at confidence level $\alpha$ is the
smallest loss threshold $\ell$ such that the probability of exceeding $\ell$
is at most $1-\alpha$:
\begin{equation}
\label{eq:var-def}
\VaR_\alpha(L) \;=\; \inf\!\bigl\{\,\ell \in \R \;:\;
  \Pr(L > \ell) \le 1-\alpha\,\bigr\}
\end{equation}
\begin{itemize}[nosep]
  \item $L$ --- portfolio loss (positive = losing money).
  \item $\alpha$ --- confidence level (e.g., 0.99 means the 99th percentile).
  \item $h$ --- holding period; always state it (``99\% 1-day VaR'').
  \item $1-\alpha$ --- the tail probability you are willing to tolerate.
\end{itemize}
Equivalently, $\VaR_\alpha$ is the $\alpha$-quantile of the loss
distribution.
\end{definition}

\begin{intuition}[What VaR Tells You]
VaR answers: ``What is the most I could lose on a normal bad day?''
The 99\% 1-day VaR says: ``On 99 out of 100 trading days, my loss will be
less than this number.''  It does \emph{not} say anything about
what happens on that 1-in-100 day---you could lose \$1 more than VaR
or \$1 billion more.  VaR is a fence, not a cliff.
\end{intuition}

\begin{workedexample}{Normal-Distribution VaR}
Suppose your portfolio's daily P\&L is normally distributed with mean
$\mu = \$50{,}000$ and standard deviation $\sigma = \$2{,}000{,}000$.

The 99\% 1-day VaR is:
\begin{align}
\VaR_{0.99} &= -({\mu - z_{0.99}\,\sigma}) \notag \\
            &= -(50{,}000 - 2.326 \times 2{,}000{,}000) \notag \\
            &= -50{,}000 + 4{,}652{,}000 \notag \\
            &= \$4{,}602{,}000
\end{align}
\begin{itemize}[nosep]
  \item We subtract $\mu$ because VaR is a \emph{loss} (sign convention:
        loss positive).
  \item The $z_{0.99}=2.326$ factor comes from the standard normal
        inverse CDF.
  \item Interpretation: on 99\% of days, the portfolio loses less than
        \$4.6 million.
\end{itemize}

The 95\% 1-day VaR uses $z_{0.95}=1.645$:
\[
  \VaR_{0.95} = -(50{,}000 - 1.645 \times 2{,}000{,}000) = \$3{,}240{,}000
\]
Lower confidence $\Rightarrow$ lower VaR.  This should be intuitive: a
less extreme ``bad day'' threshold is easier to breach.
\end{workedexample}

\begin{prereq}[Loss Conventions]
Finance has conflicting sign conventions.  In risk management, \textbf{loss
is positive}: if your portfolio drops \$5M, $L = +5{,}000{,}000$.
In portfolio theory, \textbf{returns are positive} when you make money.
VaR is reported as a positive number.  Always check the convention when
reading code or papers.
\end{prereq}


% ──────────────────────────────────────────────────────────────
\section{Expected Shortfall: VaR's Better Sibling}
\label{sec:expected-shortfall}
% ──────────────────────────────────────────────────────────────

VaR has a critical deficiency: it ignores what happens \emph{beyond} the
threshold.  Two portfolios can have identical VaR but wildly different tail
losses.  Expected Shortfall fixes this.

\begin{definition}[Expected Shortfall]
The \textbf{Expected Shortfall} (also called CVaR or Conditional VaR) at
confidence level $\alpha$ is the expected loss, conditional on the loss
exceeding VaR:
\begin{equation}
\label{eq:es-def}
\ES_\alpha(L) \;=\; \E\bigl[\,L \;\big|\; L \ge \VaR_\alpha(L)\,\bigr]
\end{equation}
For continuous distributions, equivalently:
\begin{equation}
\label{eq:es-integral}
\ES_\alpha(L) \;=\; \frac{1}{1-\alpha}\int_\alpha^1 \VaR_u(L)\,du
\end{equation}
\begin{itemize}[nosep]
  \item $\ES_\alpha$ averages \emph{all} quantiles beyond
        $\VaR_\alpha$---it captures tail severity.
  \item For the normal distribution:
        $\ES_\alpha = \mu + \sigma\,\frac{\phi(z_\alpha)}{1-\alpha}$,
        where $\phi$ is the standard normal PDF.
  \item $\ES_\alpha \ge \VaR_\alpha$ always (it looks deeper into the
        tail).
\end{itemize}
\end{definition}

\begin{keyidea}[Why Expected Shortfall Matters]
VaR fails a property called \textbf{subadditivity}: for some portfolios,
$\VaR(A+B) > \VaR(A) + \VaR(B)$.  This means VaR can penalize
diversification---combining two portfolios can \emph{increase} measured
risk, which is economically absurd.

ES is \textbf{subadditive}: $\ES(A+B) \le \ES(A) + \ES(B)$ always.
This makes it a \emph{coherent} risk measure in the sense of
Artzner et al.\ (1999).  The Basel Committee's FRTB framework
(Basel Committee, 2019) replaced VaR with ES at the 97.5\% confidence
level for internal-model capital calculations precisely for this reason.
\end{keyidea}

\begin{workedexample}{VaR Subadditivity Failure}
Consider two independent bonds, each with a 4\% probability of defaulting
and losing \$100, and a 96\% probability of no loss.

\textbf{Individual 95\% VaR:} For each bond, the 95th percentile loss is
\$0 (since 96\% of outcomes have zero loss, and $96\% > 95\%$).
So $\VaR_{0.95}(\text{Bond}_1) = \VaR_{0.95}(\text{Bond}_2) = \$0$.

\textbf{Portfolio 95\% VaR:} Holding both bonds (50/50), the probability of
at least one default is $1-(0.96)^2 = 7.84\%$.  Since $7.84\% > 5\%$, the
95th percentile loss is now positive: $\VaR_{0.95}(\text{Portfolio}) > \$0$.

Diversification \emph{increased} measured VaR---a clear failure of
subadditivity.  ES would not exhibit this pathology.
\end{workedexample}

\begin{keyresult}[Basel FRTB: ES Replaces VaR (Basel Committee, 2019)]
The Fundamental Review of the Trading Book replaced 99\% 10-day VaR
with 97.5\% Expected Shortfall for internal-model capital.  The 97.5\%
confidence level was chosen to provide ``a broadly similar level of risk
capture as the existing 99th percentile VaR threshold'' while yielding
more stable model output and less sensitivity to extreme outliers.
The Committee estimated this would increase market risk capital
requirements by approximately 22\% on average.
\end{keyresult}


% ──────────────────────────────────────────────────────────────
\section{VaR Computation Methods}
\label{sec:var-methods}
% ──────────────────────────────────────────────────────────────

There are three standard approaches to computing VaR.  Each makes different
trade-offs between accuracy, speed, and assumptions.  SecDB uses variants
of all three depending on the desk and asset class.

\begin{prereq}[Covariance Matrix]
For a portfolio with $n$ assets, the \textbf{covariance matrix}
$\bm{\Sigma} \in \R^{n \times n}$ has entries
$\Sigma_{ij} = \text{Cov}(r_i, r_j)$ where $r_i$ is the return on
asset~$i$.  The portfolio variance is
$\sigma_p^2 = \bw^\top \bm{\Sigma}\, \bw$,
where $\bw$ is the vector of portfolio weights.
Estimating $\bm{\Sigma}$ reliably is hard when $n$ is large
(Chapter~\ref{ch:regularized} discusses shrinkage estimators).
\end{prereq}

\subsection{Historical Simulation}
\label{subsec:hist-sim}

The simplest approach: replay actual historical returns through the current
portfolio.

\textbf{Algorithm:}
\begin{enumerate}[nosep]
  \item Collect the last $T$ days of returns for every position in the
        portfolio (typically $T=250$ or $T=500$).
  \item Revalue the current portfolio under each historical day's returns,
        producing $T$ hypothetical P\&L outcomes.
  \item Sort the P\&L from worst to best.
  \item The VaR at confidence $\alpha$ is the $(1-\alpha)\times T$-th
        worst outcome (e.g., for 99\% VaR with $T=500$, take the 5th-worst
        day).
\end{enumerate}

\textbf{Advantages:} No distributional assumptions; captures fat tails,
skewness, and nonlinear positions naturally; easy to explain.

\textbf{Disadvantages:} Requires history; assumes the past is representative
of the future; slow to adapt to regime changes; suffers from the cliff
effect (below).

\begin{warning}[The Historical Simulation Cliff Effect]
If the lookback window is 500 days, the 99\% VaR is literally the
5th-worst day.  Add or remove one extreme day from the window, and VaR
can jump discontinuously.  This creates \textbf{spurious features}:
VaR changes driven by old observations rolling out of the window rather
than by genuine changes in portfolio risk.

Adrian and Shin (2014) document that intermediary leverage is ``negatively
aligned with the banks' Value-at-Risk''---VaR constrains leverage in
real time.  If the VaR number itself is noisy due to cliff effects,
features derived from VaR rate-of-change will inherit that noise.

\textbf{Mitigation for your project:} When engineering VaR-dynamics
features (Chapter~\ref{ch:features}), apply smoothing (e.g., 5-day
moving average of VaR) and flag dates where extreme observations enter
or leave the lookback window.
\end{warning}

\subsection{Parametric (Variance-Covariance) Method}
\label{subsec:parametric}

Assume portfolio returns are normally distributed (or $t$-distributed),
then compute VaR analytically.

\textbf{Formula} (for a linear portfolio under normality):
\begin{equation}
\label{eq:parametric-var}
\VaR_\alpha = -\bigl(\bw^\top \bm{\mu}\bigr)
  + z_\alpha \sqrt{\bw^\top \bm{\Sigma}\, \bw}
\end{equation}
\begin{itemize}[nosep]
  \item $\bw$ --- vector of position sizes (dollar exposures).
  \item $\bm{\mu}$ --- vector of expected returns (often set to zero for
        1-day horizon).
  \item $\bm{\Sigma}$ --- covariance matrix of returns.
  \item $z_\alpha$ --- standard normal quantile at confidence $\alpha$.
\end{itemize}

\textbf{Advantages:} Fast (closed-form); decomposes naturally into
component VaR (Section~\ref{sec:component-var}); easy to compute
marginal contributions.

\textbf{Disadvantages:} Normality assumption is wrong (returns have fat
tails); cannot handle nonlinear payoffs (options) without delta-gamma
extensions; underestimates tail risk.

\subsection{Monte Carlo Simulation}
\label{subsec:monte-carlo}

Simulate thousands of scenarios from a specified model, revalue the
portfolio under each, and read off the empirical quantile.

\begin{prereq}[Monte Carlo Basics]
Monte Carlo estimation replaces an expectation
$\E[g(X)]$ with the sample average
$\frac{1}{N}\sum_{i=1}^N g(X_i)$, where $X_1,\ldots,X_N$ are
drawn from the distribution of $X$.  The error decreases as
$O(1/\sqrt{N})$---you need 100$\times$ more samples to get 10$\times$
more precision.
\end{prereq}

\textbf{Algorithm:}
\begin{enumerate}[nosep]
  \item Specify a model for risk-factor dynamics (e.g., multivariate
        normal, $t$-copula, GARCH).
  \item Draw $N$ scenarios (typically $N=10{,}000$--$100{,}000$).
  \item Revalue the full portfolio under each scenario using pricing
        models.
  \item Sort the resulting P\&L distribution and extract the
        $\alpha$-quantile.
\end{enumerate}

\textbf{Advantages:} Handles any distribution, any payoff (options,
structured products); can incorporate complex dependencies (copulas);
most flexible method.

\textbf{Disadvantages:} Computationally expensive (SecDB runs overnight
batch jobs for this); model risk---garbage in, garbage out; convergence
requires many scenarios.

\subsection{Method Comparison}

\begin{table}[h]
\centering
\small
\begin{tabular}{@{}lcccc@{}}
\toprule
\textbf{Property} & \textbf{Historical Sim} & \textbf{Parametric}
  & \textbf{Monte Carlo} \\
\midrule
Distributional assumptions & None & Normal/t & Model-dependent \\
Nonlinear payoffs (options) & Yes & No (delta only) & Yes \\
Speed & Fast & Very fast & Slow \\
Tail accuracy & Limited by data & Poor (thin tails) & Good \\
Decomposability & Difficult & Natural & Moderate \\
Regime adaptation & Slow & Moderate & Model-dependent \\
Cliff/discontinuity risk & High & None & None \\
\bottomrule
\end{tabular}
\caption{Comparison of VaR computation methods.  Adapted from
Jorion (2006, Ch.~9--12).}
\label{tab:var-methods}
\end{table}

\begin{keyidea}[Which Method Does SecDB Use?]
In practice, large dealer banks like Goldman Sachs use \textbf{Monte Carlo
simulation} for official risk numbers because their books contain
options, exotics, and structured products with nonlinear payoffs.
Historical simulation may be used as a cross-check or for simpler desks.
The parametric method is primarily useful for its analytical
decomposability---it gives you closed-form component VaR and marginal
VaR, which is why factor-VaR decompositions (Section~\ref{sec:factor-var})
often use a parametric framework even when the headline VaR is Monte Carlo.
\end{keyidea}


% ──────────────────────────────────────────────────────────────
\section{Component VaR and Marginal VaR}
\label{sec:component-var}
% ──────────────────────────────────────────────────────────────

Total VaR tells you the portfolio's aggregate risk.  But a risk manager
needs to know \emph{where} the risk lives---which positions, desks, or
asset classes are driving the number.  Component VaR and Marginal VaR
answer this.

\begin{prereq}[Euler's Theorem for Homogeneous Functions]
A function $f(\bw)$ is \textbf{homogeneous of degree $k$} if
$f(\lambda \bw) = \lambda^k f(\bw)$ for all $\lambda > 0$.
Euler's theorem states:
$\sum_i w_i \frac{\partial f}{\partial w_i} = k \cdot f(\bw)$.
Portfolio VaR (under the parametric method) is homogeneous of degree 1
in position sizes: double all positions, double VaR.  This gives the
Euler decomposition.
\end{prereq}

\begin{definition}[Marginal VaR and Component VaR]
For a portfolio with position sizes $\bw = (w_1, \ldots, w_n)^\top$:

\textbf{Marginal VaR} of position $i$:
\begin{equation}
\label{eq:mvar}
\text{MVaR}_i \;=\; \frac{\partial \VaR}{\partial w_i}
\end{equation}
This measures how much total VaR changes per unit increase in position $i$.

\textbf{Component VaR} of position $i$:
\begin{equation}
\label{eq:cvar}
\text{CVaR}_i \;=\; w_i \cdot \text{MVaR}_i
  \;=\; w_i \cdot \frac{\partial \VaR}{\partial w_i}
\end{equation}

\textbf{Euler decomposition property:} Component VaRs sum exactly to
total VaR:
\begin{equation}
\label{eq:euler}
\sum_{i=1}^n \text{CVaR}_i \;=\; \VaR
\end{equation}
\begin{itemize}[nosep]
  \item $\text{MVaR}_i$ --- the marginal effect of a small change in
        position $i$ on total VaR; identifies the best and worst hedges.
  \item $\text{CVaR}_i$ --- position $i$'s \emph{current} contribution
        to total VaR; can be negative (position is a hedge).
  \item The Euler property is exact under homogeneity of degree 1;
        it is an approximation for Monte Carlo VaR but widely used
        (McNeil, Frey, and Embrechts, 2015, Ch.~8).
\end{itemize}
\end{definition}

Under the parametric method, the formulas are explicit.  Portfolio VaR is:
\[
  \VaR = z_\alpha \,\sigma_p
  = z_\alpha \sqrt{\bw^\top \bm{\Sigma}\, \bw}
\]
Then:
\begin{equation}
\label{eq:mvar-parametric}
\text{MVaR}_i = z_\alpha \cdot
  \frac{(\bm{\Sigma}\,\bw)_i}{\sigma_p}
  = z_\alpha \cdot \frac{\text{Cov}(r_i, r_p)}{\sigma_p}
  = z_\alpha \cdot \beta_i \cdot \sigma_p / 1
\end{equation}
where $\beta_i = \text{Cov}(r_i, r_p)/\sigma_p^2$ is the beta of
asset $i$ with respect to the portfolio.

\begin{intuition}[Component VaR as a Risk Budget]
Think of component VaR as a \textbf{risk budget}.  If total VaR is
\$10M and your rates desk contributes \$7M while credit contributes
\$3M, then rates is ``using'' 70\% of the risk budget.  If rates' share
spikes to 90\%, that signals concentration---the portfolio is
increasingly a one-factor bet.  This is exactly the kind of shift you
will engineer features from in Phase~2.
\end{intuition}

\begin{workedexample}{Component VaR for a 3-Asset Portfolio}
Consider a portfolio with three assets: Equities (\$100M), Bonds
(\$80M), Commodities (\$20M).  Suppose the 99\% 1-day VaR is \$5.2M,
and we have computed (via the covariance matrix):
\[
  \text{MVaR}_{\text{Eq}} = 0.035, \quad
  \text{MVaR}_{\text{Bd}} = 0.010, \quad
  \text{MVaR}_{\text{Cm}} = 0.025
\]
\textbf{Component VaRs:}
\begin{align}
  \text{CVaR}_{\text{Eq}} &= 100 \times 0.035 = \$3.50\text{M} \notag \\
  \text{CVaR}_{\text{Bd}} &= 80  \times 0.010 = \$0.80\text{M} \notag \\
  \text{CVaR}_{\text{Cm}} &= 20  \times 0.025 = \$0.50\text{M} \notag
\end{align}
\textbf{Check:} $3.50 + 0.80 + 0.50 = \$4.80\text{M}$.

Wait---this does not equal \$5.2M.  That is because we have assumed
zero-mean returns; with the Euler decomposition, we also need the
diversification benefit.  In practice, the remaining \$0.40M comes from
the interaction (correlation) terms that we have simplified away.
In a correctly computed parametric VaR, the Euler property holds
\emph{exactly}:
$\sum_i \text{CVaR}_i = \VaR = \$5.2\text{M}$.

\textbf{Component VaR shares:}
Equities: $3.50/5.20 = 67\%$, Bonds: $0.80/5.20 = 15\%$,
Commodities: $0.50/5.20 = 10\%$, cross-terms: 8\%.

\textbf{Feature idea:} Track how these shares change over time.
A sudden spike in the equity share from 67\% to 90\% signals
that the portfolio has become a concentrated equity bet---exactly
the kind of crowding signal motivating this project
(Chapter~\ref{ch:intermediary}).
\end{workedexample}


% ──────────────────────────────────────────────────────────────
\section{Factor-VaR Decomposition}
\label{sec:factor-var}
% ──────────────────────────────────────────────────────────────

Component VaR decomposes risk by \emph{position}.  Factor-VaR
decomposes risk by \emph{source}---which underlying risk factors
(interest rates, credit spreads, equity indices, FX rates, commodity
prices) drive the portfolio's VaR.

\begin{prereq}[Factor Models for Returns]
A factor model expresses portfolio returns as a linear combination of
common factors plus idiosyncratic noise:
\[
  r_p = \sum_{k=1}^K \beta_k f_k + \varepsilon
\]
where $f_k$ is the return on factor $k$, $\beta_k$ is the portfolio's
exposure (loading) to factor $k$, and $\varepsilon$ is the residual.
This framework is developed fully in Chapter~\ref{ch:asset-pricing}.
\end{prereq}

\subsection{Factor-VaR Computation}

Given a factor model, the portfolio variance decomposes as:
\begin{equation}
\label{eq:factor-var-decomp}
\sigma_p^2 = \bbeta^\top \bm{\Sigma}_f \, \bbeta + \sigma_\varepsilon^2
\end{equation}
\begin{itemize}[nosep]
  \item $\bbeta = (\beta_1, \ldots, \beta_K)^\top$ --- factor loadings.
  \item $\bm{\Sigma}_f$ --- covariance matrix of factors.
  \item $\sigma_\varepsilon^2$ --- idiosyncratic variance (usually small
        for large portfolios).
\end{itemize}

The \textbf{factor-VaR contribution} of factor $k$ is:
\begin{equation}
\label{eq:factor-cvar}
\text{FVaR}_k = \frac{\beta_k \,(\bm{\Sigma}_f\,\bbeta)_k}
  {\sigma_p^2} \times \VaR
\end{equation}
This decomposes total VaR into the portion attributable to each factor.
The shares $s_k = \text{FVaR}_k / \VaR$ sum to approximately 1 (with the
remainder attributed to idiosyncratic risk).

\subsection{Herfindahl Index as a Concentration Metric}
\label{subsec:hhi}

Given factor-VaR shares $s_1, \ldots, s_K$, the
\textbf{Herfindahl-Hirschman Index} measures concentration:

\begin{definition}[Herfindahl-Hirschman Index on Factor-VaR]
\begin{equation}
\label{eq:hhi}
\HHI = \sum_{k=1}^K s_k^2
\end{equation}
\begin{itemize}[nosep]
  \item $s_k$ --- factor $k$'s share of total VaR
        ($s_k = \text{FVaR}_k / \VaR$).
  \item $\HHI \in [1/K,\; 1]$.
  \item $\HHI = 1/K$ when all factors contribute equally
        (maximum diversification).
  \item $\HHI = 1$ when a single factor drives all risk
        (maximum concentration).
  \item The \textbf{effective number of factors} is $N_{\text{eff}} = 1/\HHI$.
\end{itemize}
\end{definition}

\begin{keyidea}[Low HHI = Diversified; High HHI = Hidden Crowding]
A portfolio with $\HHI = 0.10$ has an effective number of 10 independent
risk sources.  A portfolio with $\HHI = 0.50$ is effectively a two-factor
bet.

Why does this matter for alpha?  He, Kelly, and Manela (2017) show that
dealer balance-sheet constraints price risk across asset classes
(Chapter~\ref{ch:intermediary}).  When many dealers crowd into the same
factor, HHI rises.  If that factor reverses, concentrated dealers must
all deleverage simultaneously---triggering fire sales
(Coval and Stafford, 2007) and predictable return reversals.

\textbf{Project implication:} Factor-VaR HHI is one of your
\emph{priority} features in Phase~2.  Rising HHI predicts crowding;
crowding predicts reversals.
\end{keyidea}

\begin{workedexample}{Herfindahl Index Calculation}
A cross-asset desk has factor-VaR shares:

\begin{center}
\small
\begin{tabular}{@{}lcc@{}}
\toprule
\textbf{Factor} & \textbf{VaR Share} & \textbf{Share$^2$} \\
\midrule
Interest Rates   & 0.45 & 0.2025 \\
Credit Spreads   & 0.25 & 0.0625 \\
Equity           & 0.15 & 0.0225 \\
FX               & 0.10 & 0.0100 \\
Commodities      & 0.05 & 0.0025 \\
\midrule
\textbf{Total}   & 1.00 & $\HHI = 0.30$ \\
\bottomrule
\end{tabular}
\end{center}

$N_{\text{eff}} = 1/0.30 = 3.33$ effective factors.  Despite
trading 5 factor classes, 70\% of risk comes from rates and credit.
The portfolio is less diversified than it appears.

Now suppose rates dominance increases to 75\%:

\begin{center}
\small
\begin{tabular}{@{}lcc@{}}
\toprule
\textbf{Factor} & \textbf{VaR Share} & \textbf{Share$^2$} \\
\midrule
Interest Rates   & 0.75 & 0.5625 \\
Credit Spreads   & 0.10 & 0.0100 \\
Equity           & 0.08 & 0.0064 \\
FX               & 0.05 & 0.0025 \\
Commodities      & 0.02 & 0.0004 \\
\midrule
\textbf{Total}   & 1.00 & $\HHI = 0.58$ \\
\bottomrule
\end{tabular}
\end{center}

$N_{\text{eff}} = 1/0.58 = 1.72$.  The portfolio is now essentially a
single-factor rates bet.  This spike in HHI---from 0.30 to 0.58---is
the kind of signal that may predict subsequent volatility or forced
repositioning.
\end{workedexample}


% ──────────────────────────────────────────────────────────────
\section{Scenario Analysis and Stress Testing}
\label{sec:scenarios}
% ──────────────────────────────────────────────────────────────

VaR describes ``normal'' bad days.  Stress testing asks: ``What happens
on abnormal days?''  The risk system computes the portfolio's P\&L under
a set of specified scenarios---large, specific market moves that may or
may not have occurred historically.

\begin{prereq}[What is a ``Scenario''?]
A \textbf{scenario} is a vector of simultaneous market moves:
$\Delta \bx = (\Delta x_1, \ldots, \Delta x_m)$, where each $\Delta x_j$
is a specified change in a risk factor (e.g., ``interest rates $+200$bp,
equities $-30\%$, credit spreads $+150$bp'').  The portfolio is revalued
under these moves, and the resulting P\&L is the \textbf{scenario P\&L}.
\end{prereq}

\subsection{Types of Scenarios}

\begin{enumerate}[nosep]
  \item \textbf{Historical scenarios:} Replay an actual crisis.  Common
        examples:
  \begin{itemize}[nosep]
    \item 2008 Lehman collapse (Sep--Oct 2008)
    \item 2020 COVID crash (Feb--Mar 2020)
    \item 2013 Taper Tantrum (May--Jun 2013)
    \item 1998 LTCM/Russia crisis
  \end{itemize}
  The advantage is realism; the disadvantage is that the next crisis will
  look different.

  \item \textbf{Hypothetical scenarios:} Designed to capture plausible
        but unprecedented events:
  \begin{itemize}[nosep]
    \item Sudden Fed rate hike of 100bp
    \item Eurozone sovereign default
    \item Oil price spike to \$200/barrel
  \end{itemize}
  These are judgment-based and harder to calibrate.

  \item \textbf{Regulatory scenarios:} Mandated by supervisors (e.g.,
        the Fed's CCAR/DFAST stress tests, which specify exact
        macroeconomic paths including GDP, unemployment, and asset
        prices).  The Fed's 2026 severely adverse scenario specifies a
        ``severe global recession triggered by an abrupt decline in risk
        appetite'' with substantial declines in risky asset prices.
\end{enumerate}

\subsection{Scenario P\&L as Features}

The raw scenario P\&L numbers are rich feature material:

\begin{keyidea}[Three Feature Families from Scenario P\&L]
\begin{enumerate}[nosep]
  \item \textbf{Scenario rank:} Which scenario produces the worst P\&L?
        If the worst scenario changes from ``rates shock'' to ``credit
        shock,'' the portfolio's vulnerability has shifted.
  \item \textbf{Scenario dispersion:} The spread between the best and
        worst scenario P\&L.  High dispersion means the portfolio is
        directionally exposed; low dispersion means it is well-hedged
        across scenarios.
  \item \textbf{Worst-case identity:} Tracking \emph{which} scenario is
        the worst over time can reveal regime shifts before they appear
        in VaR (since VaR is backward-looking and scenarios are
        forward-looking).
\end{enumerate}
These map directly to the ``Scenario P\&L'' feature family in Phase~2
(Chapter~\ref{ch:features}).
\end{keyidea}

\begin{workedexample}{Scenario P\&L Table}
A desk runs five standard stress scenarios on a given day:

\begin{center}
\small
\begin{tabular}{@{}lr@{}}
\toprule
\textbf{Scenario} & \textbf{P\&L (\$M)} \\
\midrule
Rates +200bp       & $-12.3$ \\
Rates $-100$bp     & $+5.1$  \\
Equities $-25\%$   & $-8.7$  \\
Credit +150bp      & $-15.9$ \\
FX (USD +10\%)     & $-3.2$  \\
\bottomrule
\end{tabular}
\end{center}

\textbf{Features:}
\begin{itemize}[nosep]
  \item Worst scenario: Credit $+150$bp ($-\$15.9$M).  $\Rightarrow$
        The desk is most vulnerable to credit widening.
  \item Scenario dispersion: $5.1 - (-15.9) = \$21.0$M.
  \item Scenario rank order: Credit $>$ Rates up $>$ Equities down
        $>$ FX $>$ Rates down.
  \item If yesterday the worst scenario was ``Rates $+200$bp,'' the
        \emph{change} in worst-case identity is itself a signal.
\end{itemize}
\end{workedexample}


% ──────────────────────────────────────────────────────────────
\section{VaR Utilization}
\label{sec:var-utilization}
% ──────────────────────────────────────────────────────────────

Every trading desk has a \textbf{VaR limit}---the maximum VaR it is
allowed to run.  VaR utilization is the ratio of current VaR to this
limit.

\begin{definition}[VaR Utilization]
\begin{equation}
\label{eq:var-util}
\text{VaR Utilization} = \frac{\VaR_{\text{current}}}
  {\VaR_{\text{limit}}} \times 100\%
\end{equation}
\begin{itemize}[nosep]
  \item Ranges from 0\% (no risk deployed) to 100\%+ (limit breach).
  \item Desks typically operate at 50--80\% utilization under normal
        conditions.
  \item When utilization approaches 100\%, the desk must reduce
        positions or get a limit increase.
\end{itemize}
\end{definition}

\begin{prereq}[Risk Limits in a Bank]
Banks impose \textbf{risk limits} at multiple levels: individual trader,
desk, division, and firm.  These are set by the firm's risk committee
based on capital allocation, expected returns, and regulatory
constraints.  Breaching a limit requires immediate notification to risk
management and usually forces position reduction.  Limits are the
mechanism by which balance-sheet constraints (Chapter~\ref{ch:intermediary})
become operational.
\end{prereq}

\subsection{Why VaR Utilization Matters for Alpha}

This is where risk management connects to the intermediary asset pricing
theory of Chapter~\ref{ch:intermediary}.

\begin{keyidea}[VaR Utilization as the Binding Constraint]
Adrian and Shin (2014) show that dealer leverage is procyclical and
``negatively aligned with the banks' Value-at-Risk.''  The mechanism:

\begin{enumerate}[nosep]
  \item Markets are calm $\Rightarrow$ VaR is low $\Rightarrow$
        utilization is low $\Rightarrow$ desks add risk $\Rightarrow$
        asset prices rise.
  \item A shock hits $\Rightarrow$ VaR spikes $\Rightarrow$ utilization
        hits limits $\Rightarrow$ desks must sell $\Rightarrow$ prices
        fall $\Rightarrow$ VaR rises further.
\end{enumerate}

This is the \textbf{VaR feedback loop}: risk constraints amplify shocks.
The intermediary asset pricing literature (He and Krishnamurthy, 2013;
Adrian, Etula, and Muir, 2014) shows this cycle is \emph{priced}---returns
are predictable from intermediary constraint proxies.

\textbf{The SecDB advantage:} External researchers proxy dealer
constraints with stale quarterly Fed Z.1 data.  You have daily VaR
utilization at the desk level---a much higher-frequency, more granular
signal.
\end{keyidea}

\begin{keyresult}[Procyclical Leverage (Adrian and Shin, 2014)]
Adrian and Shin document that broker-dealer leverage is ``kept constant
relative to Value-at-Risk'' even at the height of the crisis.  When VaR
rises (bad times), banks deleverage aggressively; when VaR falls (good
times), they relever.  This procyclical behavior---amplifying booms and
busts---is driven by VaR constraints at the desk level, precisely the
data you will have access to in SecDB.
\end{keyresult}

\subsection{VaR Utilization Features}

For the project, VaR utilization generates several natural features:

\begin{enumerate}[nosep]
  \item \textbf{Level:} Current utilization percentage.  High values
        ($>85\%$) signal constrained dealers.
  \item \textbf{Rate of change:} $\Delta\text{Util}_t =
        \text{Util}_t - \text{Util}_{t-5}$.  Rising utilization signals
        building stress.
  \item \textbf{Cross-desk dispersion:} If some desks are at 90\%
        utilization while others are at 30\%, the firm has uneven risk
        deployment---potential for internal capital reallocation.
  \item \textbf{Distance to limit:} $100\% - \text{Util}_t$.  Nonlinear
        effects: the last 5\% of headroom matters far more than the
        first 5\%.
\end{enumerate}

\begin{warning}[Utilization is Not Public]
VaR utilization is proprietary internal data.  It is one of the most
direct proxies for the intermediary constraint variables in the academic
literature, which is precisely why it could be valuable as a feature.
But it also means you cannot validate your features against any
published benchmark.  Your ridge baseline
(Chapter~\ref{ch:regularized}) and deflated Sharpe ratio
(Chapter~\ref{ch:validation}) are the guards against overfitting to a
unique dataset.
\end{warning}

\subsection{When Utilization Hits Limits: Fire Sales}

\begin{keyresult}[Fire Sales and Forced Selling (Coval and Stafford, 2007)]
Coval and Stafford study mutual funds forced to sell by capital outflows
(1980--2004).  Stocks subject to flow-driven forced sales experience
cumulative abnormal returns of approximately $-10\%$ during the selling
period, with reversals of $+2\%$ over the next quarter and $+8\%$ over
the following year.  An investment strategy trading against anticipated
forced sales earns abnormal returns exceeding 10\% annually.

The mechanism at dealer banks is analogous: when VaR utilization breaches
limits, the desk must sell---regardless of fundamental value.  This
creates predictable price pressure and subsequent reversals.
\end{keyresult}


% ──────────────────────────────────────────────────────────────
\section{Risk-Model Methodology Changes}
\label{sec:methodology-changes}
% ──────────────────────────────────────────────────────────────

Risk models are not static.  Banks periodically change how they compute
VaR: switching from historical simulation to Monte Carlo, changing the
lookback window length, adding or removing risk factors, recalibrating
volatility models.  These changes are driven by regulatory requirements,
model reviews, or genuine improvements---but they create problems for
anyone using VaR as a feature.

\subsection{Common Methodology Changes}

\begin{enumerate}[nosep]
  \item \textbf{Historical sim $\to$ Monte Carlo transition:} The
        headline VaR number can jump on the switchover date, not because
        risk changed, but because the \emph{measurement} changed.
  \item \textbf{Lookback window changes:} Moving from 250-day to
        500-day historical simulation changes VaR levels and dynamics.
  \item \textbf{Risk factor additions:} Adding a new risk factor
        (e.g., breaking ``credit'' into ``IG credit'' and ``HY credit'')
        changes factor-VaR decompositions even if total VaR is stable.
  \item \textbf{Volatility model updates:} Switching from equally-weighted
        to exponentially-weighted covariance estimation changes VaR
        responsiveness.
  \item \textbf{Correlation assumptions:} Changes to copula
        specifications or stress correlation overrides.
\end{enumerate}

\begin{warning}[Structural Breaks from Methodology Changes]
If VaR methodology changed during your lookback period, your VaR features
will have \textbf{structural breaks}---discontinuities that look like
signals but are actually measurement artifacts.

A LightGBM model (Chapter~\ref{ch:trees}) will happily learn to
split on these breaks, creating features that backtest beautifully but
have zero predictive power going forward.

\textbf{Mitigation:}
\begin{enumerate}[nosep]
  \item In Phase~0, interview the risk team about every methodology
        change in the past 5+ years.  Get exact dates.
  \item Create a binary indicator variable for each change date.
  \item Either: (a) include the indicator as a control variable, or
        (b) restrict your training window to a single methodology regime.
  \item Check feature importance (Chapter~\ref{ch:interpretation}):
        if a VaR feature's SHAP values concentrate around methodology
        change dates, it is likely capturing the break, not a real signal.
\end{enumerate}
\end{warning}

\begin{projectconnection}[Phase 0: Risk-Model Audit]
Task~2 of the implementation plan requires you to ``investigate
risk-model methodology changes.''  This is not optional.  You need a
table like this:

\begin{center}
\small
\begin{tabular}{@{}lll@{}}
\toprule
\textbf{Date} & \textbf{Change} & \textbf{Impact on Features} \\
\midrule
2022-Q1 & Hist.\ sim $\to$ MC for rates desk & VaR level shift \\
2023-Q3 & Lookback 250d $\to$ 500d & VaR dynamics change \\
2024-Q1 & Added inflation as separate factor & Factor-VaR shares change \\
\bottomrule
\end{tabular}
\end{center}

Every row in this table is a potential confounder for your ML models.
Document them before you build a single feature.
\end{projectconnection}


% ──────────────────────────────────────────────────────────────
\section{From Risk Outputs to Features: The Bridge}
\label{sec:risk-to-features}
% ──────────────────────────────────────────────────────────────

This chapter has covered the raw outputs of a risk system.  The bridge to
your project is the observation that these outputs, designed for risk
\emph{control}, contain information about the \emph{state of dealer
balance sheets}---which the intermediary asset pricing literature
(Chapter~\ref{ch:intermediary}) shows is priced.

\begin{keyidea}[The Five Feature Families]
Every risk output maps to a Phase~2 feature family:

\begin{center}
\small
\begin{tabular}{@{}lll@{}}
\toprule
\textbf{Risk Output} & \textbf{Feature Family} & \textbf{Key Metric} \\
\midrule
Total VaR & VaR dynamics
  & Level, rate-of-change, z-score \\
Component VaR & Factor concentration
  & HHI on factor-VaR shares \\
Scenario P\&L & Scenario P\&L
  & Rank, dispersion, worst-case identity \\
VaR / Limit & VaR utilization
  & Utilization \%, distance to limit \\
Component VaR shifts & Cross-asset flow
  & $\Delta$CVaR by asset class \\
\bottomrule
\end{tabular}
\end{center}

The detailed feature engineering is in Chapter~\ref{ch:features}.  This
chapter gives you the \emph{why} behind each feature---the risk-system
mechanics that generate the data.
\end{keyidea}


% ──────────────────────────────────────────────────────────────
\section{Summary}
\label{sec:risk-summary}
% ──────────────────────────────────────────────────────────────

\begin{table}[h]
\centering
\small
\begin{tabular}{@{}p{3.5cm}p{4cm}p{6cm}@{}}
\toprule
\textbf{Concept} & \textbf{Definition} & \textbf{Project Relevance} \\
\midrule
VaR &
  $\alpha$-quantile of loss distribution &
  Raw input; VaR dynamics (level, rate-of-change) are direct features \\[6pt]

Expected Shortfall &
  Expected loss beyond VaR &
  Coherent alternative; FRTB standard; use if available in SecDB \\[6pt]

Component VaR &
  Position $i$'s contribution to total VaR (Euler decomposition) &
  Tracks where risk lives; shifts between asset classes signal
  rebalancing \\[6pt]

Factor-VaR &
  Risk attributed to underlying factors &
  HHI on factor shares measures concentration/crowding \\[6pt]

Scenario P\&L &
  Portfolio P\&L under specified stress scenarios &
  Rank, dispersion, worst-case identity as features \\[6pt]

VaR Utilization &
  VaR / VaR limit &
  Most direct proxy for intermediary constraints; priority
  feature \\[6pt]

Methodology changes &
  VaR model updates create structural breaks &
  Must audit in Phase~0; confounder if ignored \\
\bottomrule
\end{tabular}
\caption{Chapter~3 key concepts and their project relevance.}
\label{tab:ch3-summary}
\end{table}

\vspace{1em}

\noindent\textbf{Key takeaways:}
\begin{enumerate}[nosep]
  \item VaR is the industry-standard risk measure---simple, widely used,
        but not subadditive and blind to tail severity.  Expected Shortfall
        fixes both problems.
  \item Component VaR and factor-VaR decompose risk into attributable
        pieces.  The Euler decomposition guarantees the pieces sum to the
        whole.
  \item The HHI on factor-VaR shares measures concentration.  Rising HHI
        signals crowding, which predicts reversals when constraints bind.
  \item VaR utilization is the operational expression of the balance-sheet
        constraints that the intermediary asset pricing literature proves
        are priced.
  \item Scenario P\&L provides forward-looking stress information that
        complements VaR's backward-looking nature.
  \item Methodology changes create structural breaks in VaR data.  Audit
        them before building features.
\end{enumerate}

\noindent The next chapter (Chapter~\ref{ch:microstructure}) covers
market microstructure---how trades actually execute, and why
transaction costs matter for signal viability.
